\section{An explicit positive formula for the scalar six-point range}
\label{sec:outer-six-scalar-positive}

\paragraph{Scope and attribution.}
Pak, Panova, and Swanson prove a positive marked-tree formula for
all plethysm coefficients whose \emph{output} has at most two rows.
We use their KOH/GOS identities, transpose one Gaussian factor,
and compress its possibly long unary chains.  The result below
is an explicit implementation and binary-size compilation of
that known formula.  It is not a new discovery of positivity
for two-row outputs.  Combined with PNEW-011 it evaluates every
\emph{scalar} fat-point section space for the six projective-frame
points.  It does not evaluate the general flag-bundle sections
with arbitrary tail in item 5, nor the three-row outputs arising
from stable two-row tails in item 4.

\begin{theorem}[PNEW-012: the compiled scalar formula]
For each $\mu\vdash6$, the retained file
\path{certificates/outer_six_scalar_polytopes.json} specifies a
fixed list of at most $45$ bounded rational polytopes
$P_{\mu,t}(n,r)$, each in at most $11$ variables, such that
\[
 p_{\mu,(n)}^{(6n-r,r)}
       =\sum_t\#\bigl(P_{\mu,t}(n,r)\cap\mathbb Z^{d_t}\bigr)
 \quad(n\ge0,\ 0\le r\le3n).
 \tag{PNEW-S1}
\]
All coefficient matrices are fixed integer matrices.  Every
right-hand side is an integer affine function of $(n,r)$,
and every real solution has every coordinate between $0$ and
$3n$.  In particular the descriptions and the counting witnesses
have size polynomial in the binary lengths of $n,r$.
The script \path{scripts/outer_six_scalar.py} evaluates (PNEW-S1)
and decides its support in polynomial bit complexity.

For the six-point blow-up $X_6$, set $n=c-k$, $r=c$.
If $c,k\ge0$, $c-k\ge0$, and $4c-6k\ge0$, then
\[
 \dim\operatorname{Hom}_{S_6}\!\left(
 [\mu],H^0(X_6,\mathcal O(cH-k\sum E_i))\right)
     =\sum_t\#\bigl(P_{\mu,t}(c-k,c)\cap\mathbb Z^{d_t}\bigr).
 \tag{PNEW-S2}
\]
If either of the latter two inequalities fails, that section
space is zero.  These assertions concern all nonnegative
parameters, not a bounded numerical sample.
\end{theorem}

\subsection{The two classical identities used as inputs}
Put $[a+1]_q=1+q+\cdots+q^a$ for $a\ge0$, and write
$Q_j(\rho)=\sum_x\min(j,\rho_x)$,
$m_j(\rho)=\#\{x:\rho_x=j\}$, and
$b(\rho)=\sum_x(x-1)\rho_x$.  The KOH identity is
\[
 {a+b\brack b}_q
 =\sum_{\rho\vdash b}q^{2b(\rho)}
   \prod_{j:m_j(\rho)>0}
   {j(a+2)-2Q_j(\rho)+m_j(\rho)\brack m_j(\rho)}_q.
 \tag{PNEW-S3}
\]
A summand for which a child first parameter
$j(a+2)-2Q_j(\rho)$ is negative is zero.

For $\mu\vdash6$, let $l=\ell(\mu)$.  A GOS configuration is
a sequence of partitions $\boldsymbol\nu=(\nu^0,\ldots,\nu^l)$
with
\[
 \nu^0=(1^6),\quad |\nu^i|=\sum_{j>i}\mu_j,
 \quad\nu^l=\varnothing,
 \quad P_j^i=Q_j(\nu^{i+1})-2Q_j(\nu^i)+Q_j(\nu^{i-1})\ge0
 \tag{PNEW-S4}
\]
for $1\le i<l$, $1\le j\le6$.
Set $m_0=\ell(\nu^1)$, with $m_0=0$ if $l=1$, and
\[
 \tau=\sum_{i=1}^{l-1}\sum_{j=1}^{6}
      (\nu^i)'_j\bigl((\nu^i)'_j-(\nu^{i+1})'_j\bigr).
 \tag{PNEW-S5}
\]
The GOS identity, in a form with its last Gaussian transposed, is
\[
 s_\mu(1,q,\ldots,q^n)
 =\sum_{\boldsymbol\nu:m_0\le n}
    q^\tau {n-m_0+6\brack6}_q
    \prod_{i=1}^{l-1}\prod_{j:m_j(\nu^i)>0}
                  {P_j^i+m_j(\nu^i)\brack m_j(\nu^i)}_q.
 \tag{PNEW-S6}
\]
Every summand is centered at degree $3n$.
For precision, the two primary locators in Pak--Panova--Swanson,
arXiv:2511.02312v3, are Theorem 3.1, printed p.~5, and
Theorem 3.6, printed p.~7.  The original GOS formula has lower
index $n-m_0$ in the first Gaussian; Gaussian symmetry gives
the form above.  We do not claim a direct bijection between the
original article's untransposed GOS trees and our transposed
ones.  The same polynomial identity, followed by the stated
KOH recursion, proves the version used here.

All factors except ${n-m_0+6\brack6}_q$ in (PNEW-S6) have
constant parameters.  The construction enumerates every
configuration (PNEW-S4), and every ordinary KOH expansion of
these constant factors, retaining distinct expansions as
distinct tags even when their lists of leaf labels coincide.

\subsection{Compressing the only transition that preserves the partition size}
For a KOH node of type $(a,b)$, a choice $\rho\vdash b$ has
children
\[
 \bigl(j(a+2)-2Q_j(\rho),\,m_j(\rho)\bigr)
       \quad(j\text{ increasing},\ m_j(\rho)>0).
 \tag{PNEW-S7}
\]
If $b=1$, the node is a leaf labeled $a$.
If $b\ge2$, the choice $\rho=(1^b)$ has one child,
of type $(a-2(b-1),b)$.  This is the only choice with a
child of the same partition size $b$.
Every other choice satisfies
$\sum_jm_j(\rho)=\ell(\rho)\le b-1$.

A \emph{skeleton} records only choices $\rho\ne(1^b)$.
At each internal skeleton node introduce an integer variable
$u_v\ge0$, the number of consecutive omitted all-ones choices
immediately before its recorded choice.  If its input first
parameter is $a_v$, require
\[
 \widetilde a_v=a_v-2(b_v-1)u_v\ge0,
 \qquad
 a_w=j(\widetilde a_v+2)-2Q_j(\rho_v)
 \tag{PNEW-S8}
\]
for its child $w$ on edge $j$; require every leaf label to
be nonnegative.  These inequalities also ensure every omitted
intermediate all-ones step is valid.  An integer choice of
the $u_v$ expands uniquely to a complete KOH tree, and every
complete tree is obtained uniquely by maximal-run compression.
Runs of length zero are included.

The sum of the children's sizes decreases by at least one
at every recorded internal node.  Telescoping over a skeleton
of initial size $b$ gives
\[
 \#\{\text{internal nodes}\}+\#\{\text{leaves}\}\le b.
 \tag{PNEW-S9}
\]
There are therefore only finitely many skeletons for each $b$,
independent of the first parameter $a$.
Writing $s_b$ for their number gives the exact finite recurrence
\[
 s_1=1,\qquad
 s_b=\sum_{\substack{\rho\vdash b\\\rho\ne(1^b)}}
                   \prod_{j:m_j(\rho)>0}s_{m_j(\rho)}.
 \tag{PNEW-S10}
\]
It gives $s_1,\ldots,s_6=1,1,2,4,7,15$.
The same recurrence with maxima in place of sums shows that
size-six skeletons have at most three internal nodes and
three leaves.  The script enumerates the recurrence literally,
including every ordered child specified by its distinct size $j$.

Apply this compression to the single varying factor in
(PNEW-S6), with initial type $(n-m_0,6)$.  Each leaf label
$a_i$ becomes an integer affine function of $n$ and the
at most three run variables.  Append the fixed leaf labels
from the constant factors in their fixed $(i,j)$ order,
then the varying tree's leaves in depth-first order.

\subsection{The exact polytope inequalities}
For one fixed configuration, one fixed tuple of constant
KOH trees, and one of the fifteen skeletons, let
$a_1,\ldots,a_t$ denote the resulting leaf labels.
Besides (PNEW-S8) and leaf nonnegativity, introduce mark
variables $k_1,\ldots,k_t$ and impose
\[
 \begin{split}
 &k_1=0,\qquad 0\le k_1\le k_2\le\cdots\le k_t,\\
 &k_{i+1}-k_i\le a_{i+1},\qquad
   k_{i+1}+k_i\le\sum_{j=1}^{i}a_j\quad(1\le i<t),\\
 &2k_t=2r-6n+\sum_{i=1}^{t}a_i.
 \end{split}
 \tag{PNEW-S11}
\]
Together with $u_v\ge0$ these are the complete polytope
conditions.  The factor two in the last equality enforces
the required parity on integer points, without an implicit
rounding convention.  Substitution of the affine labels gives
integer matrices and affine right-hand sides as claimed.
The JSON schema records every row as
$\texttt{lhs}\cdot(u,k)\le\alpha n+\beta r+\gamma$, or
as an equality, with $\texttt{rhs}=(\alpha,\beta,\gamma)$.

Here is a direct proof of the positive marking count.
For nonnegative integers $a,b$,
\[
 [a+1]_q[b+1]_q
   =\sum_{j=0}^{\min(a,b)}q^j[a+b-2j+1]_q.
 \tag{PNEW-S12}
\]
This follows by comparing the coefficients in the two
products of finite geometric sums.  Iterating it, after
the first $i$ factors a choice of accumulated shift $k_i$
has remaining interval length
$\sum_{j\le i}a_j-2k_i$.  Equation (PNEW-S12) permits the
next shift precisely when the two inequalities in the
second line of (PNEW-S11) hold.  Each choice occurs once.
All the KOH terms and all GOS terms are centered at $3n$.
Thus a final interval of length $L$ appears with initial
degree $3n-L/2$.  Multiplying by $1-q$, its positive
contribution is at that initial degree, while its negative
contribution is at $3n+L/2+1>3n$.
For $r\le3n$, the required positive contribution therefore
has $L=6n-2r$, which is exactly the last equality in
(PNEW-S11).

Finally the two-variable Schur identity gives
\[
 p_{\mu,(n)}^{(6n-r,r)}
           =[q^r](1-q)s_\mu(1,q,\ldots,q^n).
 \tag{PNEW-S13}
\]
Indeed $s_\mu[\operatorname{Sym}^n\mathbb Q^2]$ has
character $s_\mu(x^n,x^{n-1}y,\ldots,y^n)$; the
two-variable Weyl formula, followed by $x=1,y=q$, makes
the coefficient extraction (PNEW-S13) immediate.
Equations (PNEW-S3)--(PNEW-S13) prove (PNEW-S1).

\subsection{Zero parameters and boundedness over the reals}
We retain zero leaf labels: $[0+1]_q=1$.  The inequalities
force a zero new leaf to add zero to the accumulated mark;
if it is the first leaf, the prescribed initial mark is
already zero.  Thus unit factors introduce no extra count.
A tree of type $(0,b)$ has a unique valid expansion, choosing
$\rho=(b)$ and then the leaf zero.  To check uniqueness,
observe the exact identity
\[
 \sum_j m_j(\rho)\bigl(j(a+2)-2Q_j(\rho)\bigr)
                         =ab-4b(\rho).
 \tag{PNEW-S14}
\]
At $a=0$, nonnegative child parameters force $b(\rho)=0$,
and hence $\rho=(b)$.
We always retain the varying factor as type $(n-m_0,6)$,
including the case $n=m_0$.  This represents its value one
by the unique zero leaf and avoids an undefined empty list
of marks.  If $n<m_0$, the root inequality has no solution.
For $n=0$, only $\mu=(6)$ has a configuration with $m_0=0$;
the sole surviving point has $r=0$ and value one, as required.

The identity (PNEW-S14) follows from
\[
 \sum_jj m_j(\rho)=b,
 \qquad
 \sum_j Q_j(\rho)m_j(\rho)=b+2b(\rho).
\]
It implies that each child's product of parameters is at
most the parent's.  It remains valid for real first
parameters and real run variables satisfying (PNEW-S8).
Every run variable in the varying tree is consequently at
most $a_v/(2(b_v-1))\le3n$.

For the fixed GOS part, centeredness is the exact integer
identity
\[
 2\tau+\sum_{i,j}P_j^i m_j(\nu^i)=6m_0.
 \tag{PNEW-S15}
\]
The quantity $\tau$ is nonnegative: with
$\alpha_i=(\nu^i)'$ and $\alpha_l=0$,
\[
 \tau=\tfrac12\|\alpha_1\|^2+
             \tfrac12\sum_{i=1}^{l-1}\|\alpha_i-\alpha_{i+1}\|^2.
\]
Using (PNEW-S14) on every tree, including the real runs,
gives $\sum_i a_i\le6n$.
The last equality of (PNEW-S11) therefore gives
$k_t\le r\le3n$; monotonicity bounds all other marks.
All real coordinates are thus in the compact box
$[0,3n]^{d_t}$, proving boundedness of the rational polytopes,
not merely finiteness of their integer-point sets.

\subsection{The exact finite size of the compiled list}
The first two numeric columns below enumerate (PNEW-S4)
and the constant-factor KOH expansions.  Multiplication by
$s_6=15$ gives the number of polytopes.  Different tags
are retained even if two matrices agree, since they represent
different summands of the positive identity.
\begin{center}
\begin{tabular}{c r r r r}
\toprule
$\mu$ & Configurations & Constant decorations & Polytopes & Maximum variables\\
\midrule
$(6)$             &1&1&15&6\\
$(5,1)$           &1&1&15&7\\
$(4,2)$           &2&3&45&7\\
$(4,1,1)$         &1&2&30&8\\
$(3,3)$           &3&3&45&8\\
$(3,2,1)$         &2&2&30&9\\
$(3,1,1,1)$       &1&2&30&10\\
$(2,2,2)$         &3&3&45&9\\
$(2,2,1,1)$       &2&3&45&10\\
$(2,1,1,1,1)$     &1&1&15&10\\
$(1,1,1,1,1,1)$   &1&1&15&11\\
\bottomrule
\end{tabular}
\end{center}
These small finite lists are obtained by exhaustive partition
recursion, with no unknown multiplicities as inputs.
The retained script regenerates every row and all $330$
full integer matrix descriptions.
At most three run variables and eight leaf marks occur.

\subsection{Polynomial bit complexity and an exact evaluator}
The fixed list and the coordinate bound prove the claimed
polynomial-size counting description directly.  For efficient
numerical evaluation, the script uses the same coefficient
(PNEW-S13), evaluated by the hook-content identity
\[
 s_\mu(1,q,\ldots,q^n)
    =q^{b(\mu)}
       \frac{\prod_{(i,j)\in\mu}(1-q^{n+1+j-i})}
            {\prod_{(i,j)\in\mu}(1-q^{h_{ij}})}.
 \tag{PNEW-S16}
\]
If $\ell(\mu)>n+1$ the result is zero immediately.  Otherwise
all numerator exponents are positive.  There are six numerator
factors, so at most $64$ shifted reciprocal-denominator
coefficients, or $128$ after multiplication by $1-q$, suffice.
The denominator degree is at most $21$.  Indeed the sum of
the six hook lengths is six plus the number of pairs of
cells sharing a row or a column, at most $6+\binom62$.

For completeness, the implemented coefficient extraction is
the following elementary parity recurrence.  To find
$[q^N]P(q)/Q(q)$ with $Q(0)=1$, form
$P(q)Q(-q)$ and $Q(q)Q(-q)$, retain in the former the terms
of parity $N\bmod2$, and retain in the latter its even
terms.  Divide all retained exponents by two (after removing
one from the former odd exponents).  This replaces the
coefficient problem by one at $\lfloor N/2\rfloor$.
It follows by separating the even and odd parts of the
power series, so repeated application is exact.
Negative requested exponents contribute zero.

The degrees remain bounded by $21$ and there are
$O(\log(N+1))$ iterations.  The denominator roots are roots
of unity and squaring them preserves that property; with
fixed degree, its coefficients remain bounded by a constant.
The numerator coefficients grow by at most a fixed factor
at each iteration, so their bit lengths are
$O(\log(N+1))$.  Ordinary exact integer arithmetic therefore
proves polynomial bit complexity.  Subset shifts and all
endpoint arithmetic have bit length $O(\log(n+r+1))$.
No sign assumption on intermediate hook-content sums is
needed: positivity is independently proved by (PNEW-S1).
Testing the resulting exact integer for zero decides support
throughout this scalar/two-row range.

Equation (PNEW-S2) now follows from PNEW-011, since
$6(c-k)-c=5c-6k$ and $c\le3(c-k)$ is equivalent to
$4c-6k\ge0$.  This completes the proof of PNEW-012.

\subsection{Certificates and their exact scope}
The following computations were executed with integers and
rational numbers only.
\begin{itemize}
\item All $2,717$ coefficients with $\mu\vdash6$,
$0\le n\le12$, $0\le r\le3n$ agree between positive
compressed-tree/Clebsch--Gordan path counting and the
independent hook-content parity recurrence.
The complete vectors are in
\path{certificates/outer_six_scalar_small.json}.
\item For $0\le n\le3$, all $242$ two-row cases additionally
agree with direct integer-point counting from the exported
matrix inequalities.  This check does not use the
Clebsch--Gordan count to evaluate those matrices.
\item All $5,214$ coefficients for every $\mu\vdash6$ and
every output partition of $6n$, $0\le n\le3$, were computed
by power-sum plethysm followed by beta-set
Murnaghan--Nakayama.  For every output of length at most six,
the result agrees with literal monomial cycle characters
and the $720$-term $GL_6$ Weyl alternant.  Longer outputs
are zero by the row-Pieri length bound and also vanish in
the first exact computation.  There are $44$ additional
whole-representation dimension checks.
The complete rows are in
\path{certificates/outer_six_all_n_le_3.json}; the independent
script is \path{scripts/outer_six_character_check.py}.
\end{itemize}
The full character table is a calibration dataset for
future flag-section formulas.  Agreement of the scalar
formula is claimed only on its two-row output subfamily;
the other rows are not falsely reported as checks of a
positive all-output formula.

\paragraph{Why the two character computations are independent.}
The first uses
$p_d[h_n]=\sum_{\rho\vdash n}p_{d\rho}/z_\rho$,
the outer $S_6$ character table, and rim-hook characters.
The second constructs the actual weights of
$\operatorname{Sym}^n\mathbb Q^6$, traces the eleven cycle
types on their sixth tensor power, and takes the Weyl
alternant.  The shared outer character table is separately
checked by exact character orthogonality and the hook
dimension formula.  The $44$ dimension checks compare
$\sum_\lambda p_{\mu,(n)}^\lambda\dim S_\lambda\mathbb Q^6$
with $\dim S_\mu(\mathbb Q^{\binom{n+5}{5}})$.

\paragraph{References and source review.}
Pak, I., Panova, G., \& Swanson, J. P. (2026).
\emph{A combinatorial interpretation for certain plethysm
and Kronecker coefficients} (arXiv:2511.02312v3).
\url{https://arxiv.org/pdf/2511.02312v3}.
The exact locators are Lemma 1.3 and its proof in \S2.1;
Theorem 3.1 and Definition 3.2; Proposition 3.4;
Definition 3.5 and Theorem 3.6; Proposition 3.9;
Lemmas 4.1--4.2 with their whole proofs; and Theorem 5.2
with its proof.  The entire required chain was read from
the attached LaTeX source and checked against the primary
version.  The classical KOH and GOS identities are cited
inputs; the compression, all real boundedness assertions,
the exact matrix construction, and the binary-complexity
argument are proved above.
